\section{The Sacraments of Healing: Christ the Physician of Soul and Body}
\label{sec:healing-introduction}

When the Scribes murmured because Jesus sat at table with publicans and sinners, the Lord answered with the title that gathers the whole economy of mercy into a single image: ``They that are whole, need not the physician: but they that are sick. I came not to call the just, but sinners to penance'' (Luke~5:31--32). Throughout the Gospels, the Eternal Word made flesh acts as the Divine Physician (\latin{Medicus caelestis}) of the whole human person. He refuses to sever the malady of the spirit from the affliction of the body. When four men lowered a paralytic through the roof at Capharnaum, Jesus looked first to the deeper wound and said, ``Son, thy sins are forgiven thee,'' and then confirmed His invisible power by healing the paralyzed limbs: ``Arise, take up thy bed, and go into thy house'' (Mark~2:5, 11). The miracles of Christ reveal that sin and bodily decay belong together in the ruin wrought by the Fall, and that the salvation won upon the Cross reaches down into both.

\subsection{Baptized Life in a Mortal Body and a Wounded World}
\label{subsec:baptized-life-wounded-world}

Why do the baptized---already washed in the font, sealed with the Holy Spirit, and nourished by the Bread of Life---still stand in need of medicinal sacraments? St.~Thomas Aquinas answers by comparing the life of grace to bodily life (\emph{ST}~III, q.~65, a.~1). In natural existence, birth, growth, and daily nourishment would suffice if the human body were already impassible and immortal. Because man remains vulnerable to sickness and injury, however, nature requires not only the operations that generate and sustain life directly (\latin{per se}), but also the remedial arts that expel disease and restore vigor when harm intervenes (\latin{per accidens}).

So it is in the supernatural order during earthly pilgrimage. Christian Initiation lays the entire foundation of divine life: Baptism generates the believer anew in Christ, Confirmation strengthens him in the Spirit, and the Eucharist feeds him with Christ Himself. Yet the newly regenerated Christian carries that heavenly treasure in an earthen vessel (2~Cor~4:7). As the Council of Trent teaches in its decree on Original Sin (Session~V), Baptism genuinely removes all guilt of original sin and every previous actual sin, leaving nothing in the reborn soul that God hates; nevertheless, there remains in the baptized the disordered inclination called concupiscence---the tinder of sin (\latin{fomes peccati})---left for the combat of virtue, together with the passibility, suffering, and mortality of the body. Human free will, though healed and elevated by grace, remains mutable in this life and capable of turning away from God through mortal sin (\emph{ST}~III, q.~84, a.~10). Even when charity abides, daily venial faults, lingering habits of weakness, and the crushing trials of grave sickness and approaching death assail the pilgrim soul.

Against these twin afflictions---spiritual death or debility through post-baptismal sin, and bodily infirmity burdened by the remnants of sin---Christ instituted the two \textbf{Sacraments of Healing}: Penance and the Anointing of the Sick. St.~Thomas traces their complementary design with characteristic precision:
\begin{quotation}
And this would be enough for man if he had an impassible life, both corporally and spiritually; but since man is liable at times to both corporal and spiritual infirmity, i.e.\ sin, hence man needs a cure from his infirmity; which cure is twofold. One is the healing, that restores health: and corresponding to this in the spiritual life there is Penance, according to Ps.~40:5: ``Heal my soul, for I have sinned against Thee.'' The other is the restoration of former vigor by means of suitable diet and exercise: and corresponding to this in the spiritual life there is Extreme Unction, which removes the remainder of sin, and prepares man for final glory. (\emph{ST}~III, q.~65, a.~1)
\end{quotation}
In Aquinas's ordering of the seven sacraments, Penance and Anointing stand to each other much as Baptism stands to Confirmation (\emph{ST}~III, q.~65, a.~3): Penance is more necessary, because without it the Christian who has committed mortal sin cannot recover spiritual life, whereas Anointing of the Sick brings medicinal healing to completion, fortifying the gravely ill believer in his infirmity and purging the lingering weaknesses left behind by sin.

\subsection{From Woundedness to Eucharistic Communion and Glory}
\label{subsec:healing-to-eucharist-glory}

Both Sacraments of Healing share three structural marks that set them apart within the sevenfold sacramental organism. First, neither imprints an indelible character upon the soul (\emph{ST}~III, q.~63, a.~6). Character consecrates a member permanently to a station in divine worship, whereas healing restores a wounded member to soundness; because human frailty can fall anew into sin and because grave sickness can recur or deepen across a lifetime, both Penance and Anointing are repeatable gifts of Christ's unwearied compassion (\emph{ST}~III, q.~84, a.~10). Second, both are entrusted exclusively to validly ordained bishops and priests, since each touches the remission of post-baptismal sin and exercises Christ's pastoral governance over His wounded flock (John~20:21--23; Jas~5:14--15; Trent, Sess.~XIV). Third, both are inwardly oriented toward the Holy Eucharist and toward final glory.

\begin{threestable}{Sacrament of Healing}{Wound Addressed}{Restorative Term}
Penance and Reconciliation & Post-baptismal mortal and venial sin; rupture of charity with God and visible wound to ecclesial communion & Restores sanctifying grace and peace with the Church, reopening lawful, fruitful access to the Eucharistic altar\\
Anointing of the Sick & Grave bodily illness, debility of old age, spiritual anguish, temptation in frailty, and the remnants of sin & Conforms suffering to Christ's Passion, strengthens soul and body as expedient, and prepares the believer for Viaticum and glory\\
\end{threestable}

Healing in Christ never terminates in mere self-restoration. Sin isolates the believer from the table of the Lord and wounds the whole Mystical Body; grave sickness threatens to turn bodily suffering into terror, despair, or isolation. Through Penance, the Divine Physician absolves the contrite sinner, restores peace with the Church (\latin{pax cum Ecclesia}), and leads him back to the banquet of Christ's Body and Blood. Through the Anointing of the Sick, He visits the believer weighed down by serious illness or advancing years, pours the oil of gladness and fortitude upon forehead and hands, unites human frailty to His redeeming Cross, and prepares the pilgrim either for renewed earthly service or for Viaticum---the Eucharistic food for the passage into the Father's house (\emph{ST}~III, q.~65, a.~3; \emph{Lumen Gentium}~11).
